\let\negmedspace\undefined
\let\negthickspace\undefined
\documentclass[journal]{IEEEtran}
\usepackage[a5paper, margin=10mm, onecolumn]{geometry}
\usepackage{tfrupee} 
\setlength{\headheight}{1cm} 
\setlength{\headsep}{0mm}     
\usepackage{gvv-book}
\usepackage{gvv}
\usepackage{cite}
\usepackage{amsmath,amssymb,amsfonts,amsthm}
\usepackage{algorithmic}
\usepackage{graphicx}
\usepackage{textcomp}
\usepackage{xcolor}
\usepackage{txfonts}
\usepackage{listings}
\usepackage{enumitem}
\usepackage{mathtools}
\usepackage{gensymb}
\usepackage{comment}
\usepackage[breaklinks=true]{hyperref}
\usepackage{tkz-euclide} 
\def\inputGnumericTable{}                                 
\usepackage[latin1]{inputenc}                                
\usepackage{color}                                            
\usepackage{array}                                            
\usepackage{longtable}                                       
\usepackage{calc}                                             
\usepackage{multirow}                                         
\usepackage{hhline}                                           
\usepackage{ifthen}                                           
\usepackage{lscape}
\usepackage{booktabs}
\usepackage{tikz}
\usetikzlibrary{arrows.meta,angles,quotes}

\begin{document}

\bibliographystyle{IEEEtran}
\vspace{3cm}

\title{4.12.13}
\author{AI25BTECH11008 - Chiruvella Harshith Sharan}
{\let\newpage\relax\maketitle}

\renewcommand{\thefigure}{\theenumi}
\renewcommand{\thetable}{\theenumi}
\setlength{\intextsep}{10pt} 

\numberwithin{equation}{enumi}
\numberwithin{figure}{enumi}
\renewcommand{\thetable}{\theenumi}

\textbf{Question:} \\

If the equation of the base of an equilateral triangle is $x+y=2$ and the vertex is $(2,-1)$, then find the length of the side of the triangle. Also find the other vertices.

\vspace{0.3cm}
\textbf{Answer:}

\subsection*{\textbf{Step 1: Represent line in matrix form}}

The line equation:
\[
x+y=2 \quad \implies \quad \vec{n}^T\vec{x} = 2,
\]
where
\[
\vec{n} = \begin{pmatrix}1 \\ 1\end{pmatrix}.
\]

Vertex:
\[
\vec{A} = \begin{pmatrix}2 \\ -1\end{pmatrix}.
\]

\subsection*{\textbf{Step 2: Foot of perpendicular (matrix equation)}}

The foot $\vec{D}$ is obtained by projecting $\vec{A}$ onto the line:

\[
\vec{D} = \vec{A} - \frac{\vec{n}^T\vec{A}-2}{\vec{n}^T\vec{n}}\vec{n}.
\]

Substitute:
\[
\vec{D} = \begin{pmatrix}2 \\ -1\end{pmatrix} - 
\frac{(1)(2)+(1)(-1)-2}{1^2+1^2}\begin{pmatrix}1 \\ 1\end{pmatrix}
= \begin{pmatrix}\tfrac{5}{2} \\ -\tfrac{1}{2}\end{pmatrix}.
\]

\subsection*{\textbf{Step 3: Side length of equilateral triangle}}

In an equilateral triangle, altitude $AD = \tfrac{\sqrt{3}}{2}a$.  

Compute:
\[
AD = \|\vec{A}-\vec{D}\| = \sqrt{\left(2-\tfrac{5}{2}\right)^2 + \left(-1+\tfrac{1}{2}\right)^2} 
= \frac{1}{\sqrt{2}}.
\]

So:
\[
\frac{1}{\sqrt{2}} = \frac{\sqrt{3}}{2}a \quad \implies \quad a = \sqrt{\tfrac{2}{3}}.
\]

\subsection*{\textbf{Step 4: Coordinates of other vertices}}

Let $\vec{m}$ be a direction vector along the base line. Since $\vec{n}=(1,1)^T$, take
\[
\vec{m} = \begin{pmatrix}1 \\ -1\end{pmatrix}.
\]

Half the base length is $\tfrac{a}{2}$.  
So, the base vertices are:
\[
\vec{B} = \vec{D} + \frac{a}{2}\frac{\vec{m}}{\|\vec{m}\|}, \quad 
\vec{C} = \vec{D} - \frac{a}{2}\frac{\vec{m}}{\|\vec{m}\|}.
\]

Since $\|\vec{m}\| = \sqrt{2}$:
\[
\vec{B} = \begin{pmatrix}\tfrac{5}{2}\\ -\tfrac{1}{2}\end{pmatrix} + 
\frac{\sqrt{2/3}}{2\sqrt{2}}\begin{pmatrix}1\\-1\end{pmatrix}
= \begin{pmatrix}\tfrac{5}{2}+\tfrac{1}{2\sqrt{3}}\\ -\tfrac{1}{2}-\tfrac{1}{2\sqrt{3}}\end{pmatrix},
\]

\[
\vec{C} = \begin{pmatrix}\tfrac{5}{2}\\ -\tfrac{1}{2}\end{pmatrix} - 
\frac{1}{2\sqrt{3}}\begin{pmatrix}1\\-1\end{pmatrix}
= \begin{pmatrix}\tfrac{5}{2}-\tfrac{1}{2\sqrt{3}}\\ -\tfrac{1}{2}+\tfrac{1}{2\sqrt{3}}\end{pmatrix}.
\]

\subsection*{Final Answer:}
\[
\text{Side length } a = \sqrt{\tfrac{2}{3}}, 
\quad B = \left(\tfrac{5}{2}+\tfrac{1}{2\sqrt{3}}, -\tfrac{1}{2}-\tfrac{1}{2\sqrt{3}}\right), 
\quad C = \left(\tfrac{5}{2}-\tfrac{1}{2\sqrt{3}}, -\tfrac{1}{2}+\tfrac{1}{2\sqrt{3}}\right).
\]

\begin{figure}[htbp]
\centering
\begin{tikzpicture}[scale=2]
\coordinate (A) at (2,-1);
\coordinate (D) at (2.5,-0.5);
\coordinate (B) at (2.5+0.2887,-0.5-0.2887); % approx numeric
\coordinate (C) at (2.5-0.2887,-0.5+0.2887);

\draw[thick] (-1,3) -- (3,-1) node[right]{$x+y=2$};
\draw[thick] (A)--(B)--(C)--cycle;

\fill (A) circle (1pt) node[below right] {$A(2,-1)$};
\fill (B) circle (1pt) node[right] {$B$};
\fill (C) circle (1pt) node[left] {$C$};
\fill (D) circle (1pt) node[above right] {$D$};
\end{tikzpicture}
\caption{Equilateral triangle with base on $x+y=2$}
\end{figure}

\end{document}
